\section{The preliminary comparison and a derived-code task}
\label{app:prelim}

\subsection{Preliminary comparison}
\label{sec:prelim}

Experiment 1 was motivated by an earlier, simpler comparison with different sentence templates. It contrasted superseded assignments with a single control in which the two earlier values were mentioned as unassigned values, averaging the two mention orders. Superseded relevance exceeded this control in all three models that passed that comparison's competence screen: by 3.50 nats (95\% CI [3.20, 3.80]) in Qwen3-8B, 1.76 [1.51, 2.03] in Gemma 3 4B, and 5.53 [5.27, 5.79] in Phi-4-mini. Read alone, this looks like evidence that superseded assignments retain query-specific influence. But the control matched neither the entity association nor the relative position of the superseded mentions, so the comparison could not say which property produced the difference. Experiment 1 adds controls for each.

The preliminary comparison used different templates from Experiments 1 and 2. A superseded prompt for the Nora/Liam example read:

\begin{prompt}
The badge assigned to Nora was amber.
The badge assigned to Liam was coral.
Later, Nora’s badge was changed to jade.
Later, Liam’s badge was changed to pearl.
What is Nora’s current badge?
Respond with only the value, with no explanation.
\end{prompt}

The control stated the current assignments and mentioned \lit{amber} and \lit{coral} as unassigned badge values, in both mention orders; the primary contrast averaged the two orders within each history. A live control edited the current assignment. Scoring, $E$, and $R$ were defined as in Section~\ref{sec:task}. Each history contributed 40 scored prompts, giving 3,840 prompts and 1,920 matched edit pairs over 96 confirmation histories per model, with the same history-level bootstrap. Four models were screened with the same 99\% rule. Qwen3-8B, Gemma 3 4B, and Phi-4-mini passed; Mistral-7B-Instruct-v0.3 scored 281/288 on the unassigned control and was not run.

\subsection{Derived-code task}
\label{app:derived-code}

To ask whether sensitivity to an earlier value reaches a computation that uses it, we mapped values to shuffled opaque codes and asked which code corresponded to an entity's current badge. The codebook stayed fixed across baseline and edit; editing an earlier badge value changed which code it mapped to, not the codebook or the current assignment. We scored all 16 full code sequences by their unnormalized log probability. Only Qwen passed this task's screen (48/48); its threshold, 47/48 (97.9\%), was specific to this task and recorded before scoring. Gemma scored 45/48 and Phi 33/48. For Qwen, query relevance was 5.544 nats (95\% CI [5.159, 5.941]) for an edit to the earlier value and 12.836 [12.418, 13.257] for an edit to the current value, both positive in all 96 histories. Qwen's candidate ranking is below ceiling here: with the earlier value edited, the current code ranked first in 360/384 baseline and 364/384 edited prompts; 9 pairs changed from incorrect to correct and 5 from correct to incorrect. Changes in the current code's log probability (0.108 [−0.050, 0.264]) and in the candidate margin (0.103 [−0.118, 0.325]) had intervals spanning zero. This task lacks the entity-mention and order controls of Experiment 1 and collected no generated answers, so it shows score sensitivity in a derived computation, below ceiling in candidate ranking, without saying which property of the earlier mention produces it or whether generated answers change.

\FloatBarrier
